We presented a nested EV charging controller with hourly pricing,
15-min plan-residual dispatch and 60-s reactive voltage correction. On two IEEE feeders
with AC power flow and real price, load and session data, it kept the voltage
floor without base-load deviation, up to \GPenSFive\% penetration on the
33-bus feeder and at \GPenSThree\% in the other settings,% claim: nested_zero_viol_nominal
% claim: general_nested_zero_viol
\ and was the cheapest learned controller.% claim: nested_cheapest_learned_main
\ A price-aware heuristic remained cheaper at nominal load,% claim: pa_cheaper_than_nested_all_main
\ and a perfect-foresight LP-OPF cheaper still.% claim: lp_cheapest_S3
\ Under stress beyond the chargers' reactive capability, the controller
curtailed and left undelivered far less energy than the heuristic;% claim: nested_less_curtailment_than_pa_S7
% claim: nested_less_unmet_than_pa_S7
\ the lower curtailment came from the learned levels, the lower unmet energy
from least-laxity-first allocation.% claim: decomp_nested_less_curt_than_plan_S7
% claim: plan_less_unmet_than_pa_S7
\ Two lessons follow. First, a physics-based correction layer bounded by the
inverter rating removed the violations of every coordinator we attached it to,
up to \GPenSFive\% penetration without base-load deviation, at an
energy-cost change of at most \GLThreeCostChangeMax\%;% claim: l3_negligible_cost_main
% claim: l3_zero_viol_rules_nominal
% claim: l3_zero_viol_learned_nominal
\ such a layer belongs in every comparison. Second, learned dispatch is worth
its complexity mainly when it starts from a sound plan and grid stress, not
price, is the binding concern. Future work includes unbalanced three-phase
feeders, uncertain departure times and field validation.
